\subsection{模型假设：分层介质与有限信息}

本节先把光路放回实物：空气、外延层和衬底构成相邻的分层介质；假设的职责是规定波矢穿过界面、返回场保留阶数，以及缺失光学量由何种参数承接。本文采用非磁、标量、平面平行且层内均匀的介质近似。三层界面两侧沿界面方向的波矢分量连续，两种偏振采用同一套界面规则。导纳可理解为界面上电场与磁场关系的比例量；纵向传播因子把折射率与层内传播方向合在一起，透明时等于该层折射率乘以光线与垂直界面方向夹角的余弦。$s$偏振的电场垂直入射面，其导纳取纵向传播因子；$p$偏振的电场平行入射面，其导纳取该层折射率平方与纵向传播因子的比值，二者共同进入界面系数。空气折射率取$1$，外延层和衬底折射率仍按随波数变化的未知函数处理。表\ref{tab:02}把假设绑定到公式位置和核验方式。
% 依据:交接/建模笔记_问题1.md:1.题面边界与数据事实;2.几何、波数和传播因子;3.界面系数与一次往返场
% src:交接/假设台账_问题1.json:0.假设
\begingroup
\zihao{5}
\setlength{\tabcolsep}{3pt}
\begin{longtable}{>{\centering\arraybackslash}p{0.26\textwidth} >{\centering\arraybackslash}p{0.33\textwidth} >{\centering\arraybackslash}p{0.33\textwidth}}
\caption{模型假设的作用位置与核验方式}\label{tab:02}\\
\hline
假设 & 作用位置 & 核验方式\\
\hline
\endfirsthead
\hline
假设 & 作用位置 & 核验方式\\
\hline
\endhead
平面平行层 & 传播、界面系数 & 场表达式互核\\
被动传播 & 传播因子 & 衰减与单位回换\\
偏振加权 & 导纳、反射率 & 权重重估\\
经验色散 & 往返相位 & 常数与扰动对照\\
有效损耗 & 传播因子模长 & 损耗与收敛核对\\
逐问返回阶数 & 两束场、往返场 & 级数与闭式核对\\
允许观测偏差 & 灵敏度检验 & 保留原值重估\\
\hline
\end{longtable}
\endgroup

被动传播要求幅度沿光路不增长；偏振等权是加权模型中的一个情景。有效损耗合并吸收与界面造成的振幅衰减，避免把它拆成资料无法区分的来源。

场表达式互核比较直接计算的光场强度与交叉项展开式，核对两种写法的代数一致性；结构适用性则由生成光路与估计光路不同的失配情景检查。
% src:交接/假设台账_问题1.json:0.假设;0.灵敏度义务

界面规则确定后，返回阶数按题面逐问收紧。问题一把到达衬底界面后返回的首束记为$B$，把直接反射记为$r_{01}$，两束截断场记为$E_2$；问题三才允许后续每次往返再乘同一个往返复倍率$u$。这里的往返复倍率是光再绕一周时同时改变振幅和相位的因子，完整往返场记为$E_\infty$，其关系写成
\begin{equation}
E_2=r_{01}+B,\qquad E_\infty=r_{01}+B(1+u+u^2+\cdots)=r_{01}+\frac{B}{1-u},\qquad |u|<1 .
\end{equation}
其中$|u|<1$表示每增加一次往返，返回场的振幅就按小于一的倍率递减，需逐点检查，不由材料名称预设。该条件保证返回场级数收敛；高阶条纹能否观测，还取决于光束相干、空间重叠及测量分辨条件，问题三分别检查。具体而言，各束光须保持稳定的相位关系并落在同一探测区域，高阶条纹的变化还须能与噪声区分且能被仪器分辨。返回阶数固定后，色散承接相位变化，损耗承接传播振幅的衰减。
% 依据:交接/建模笔记_问题1.md:3.界面系数与一次往返场
% 依据:交接/建模笔记_问题3.md:2.从界面场到多次往返场;4.多光束必要条件及精度影响
% src:交接/假设台账_问题3.json:5.假设

损耗来源的处理来自一次参数化取舍：曾比较分项描述消光、界面损耗和偏振权重，与合并有效往返损耗两种方案。附件缺少分辨各损耗来源的独立观测，因而采用合并结构，折射率随波数的变化仍用少量参数描述。原始反射率允许存在测量偏差，异常值、选段和残差用于检查假设对厚度估计及光谱预测的影响。
% src:交接/实验记录.json:43.尝试;43.现象;43.决定
% src:交接/数据档案.json:未随附件提供的建模输入
% src:交接/假设台账_问题3.json:2.假设;3.假设;4.假设

这些约束把“采用什么假设”落实为“假设在哪一步生效、如何被扰动核对”；下一节据此统一外角、复折射率、相位因子以及厘米和微米的符号与量纲。
